%!TEX root = ../main.tex
% Declaración sobre el uso de inteligencia artificial.
% Exigencia de la cátedra (clase 5): el uso de IA está permitido y
% fomentado, pero debe declararse. Esta página resume el uso; el
% registro completo de prompts va en el apéndice D.

\begin{guia}[Guía de la cátedra: política de uso de IA]
La IA es parte del estado del arte de la industria y su uso está
permitido, con condiciones:
\begin{enumerate}
    \item \textbf{Declararlo}: acá y en la Metodología (qué tareas fueron
          asistidas y con qué herramientas).
    \item \textbf{Autoría}: las ideas, decisiones y conclusiones son del
          estudiante. La IA asiste; no decide.
    \item \textbf{Verificación}: toda cifra, cita o referencia sugerida por
          una IA se confirma en la fuente original antes de usarla.
    \item \textbf{Trazabilidad}: los prompts relevantes se guardan en una
          carpeta compartida con la cátedra.
    \item \textbf{Comprensión}: no se incorpora texto que no se haya leído
          y entendido. Plagio, incluso si viene de un modelo de lenguaje,
          es plagio.
\end{enumerate}
\end{guia}

% BORRADOR: el autor debe revisar y confirmar cada afirmación.
Se reconoce el uso de Claude (Anthropic), a través de la herramienta
Claude Code, como asistente durante el desarrollo del proyecto. Las
tareas asistidas fueron la programación y las pruebas del prototipo, la
revisión de la consistencia entre el documento y el código, la redacción
de borradores de los capítulos, que luego fueron revisados y reescritos
por el autor, y la conversión del documento al formato de la cátedra. Los
cambios de código asistidos quedan identificados en el historial del
repositorio mediante la línea \texttt{Co-Authored-By}. \completar{Otras
herramientas utilizadas y tareas asistidas.}

Todas las salidas fueron revisadas y verificadas por el autor, que asume
la responsabilidad total sobre el contenido de este documento. Las cifras
y referencias se contrastaron con la ejecución del sistema y con las
fuentes originales. Los prompts relevantes se registran en
\completar{carpeta compartida con la cátedra}.
